% Drawing 133/03-1 -- folder 133, batch 03, archive page 45 (PDF page 46).
% Transcription: schéma en haut à droite de la page, sans flèches : les deux arcs extérieurs, marqués $M'$ et $M$, joignent $N$ et $N'$ à $G$ ($M' = G/N$, $M = G/N'$). « $NN' = H$ » est écrit sur un « $MM'$ » effacé. Le « 3 » bouclé $\mathfrak{z}$ désigne ici $N \cap N' = \pi_1$.
% Kind: lattice diagram of subgroups. Method 2 (schematic redraw from the facsimile). Status: ok.
% Notes: Subgroup diagram without arrows: N and N' under N cap N' = z (looped 3) = pi_1 at the top, both joined to NN' = H in the middle, itself joined to G at the bottom; the two outer arcs, labelled M' and M, join N and N' to G. The erased MM' under « NN' = H » is left out.
% Coordinates: 1 unit = 1 mm = 1/100 of the crop width (bbox in page fractions: [0.6, 0.065, 0.79, 0.185]).
% Redrawn by Claude (claude-opus-5-5), issue #73, 2026-11-07; not checked by a human; not compiled (no TeX in the session).
% Rights: reproduces part of the work -- behind RENDER_TRANSCRIPTS, never published by itself.
\documentclass[tikz,border=3pt]{standalone}
\usepackage{amsmath,amssymb}
\usetikzlibrary{patterns,arrows.meta}
\tikzset{pencil/.style={black!45}, struck/.style={black!35}, mark/.style={circle,draw,inner sep=1.3pt}, pt/.style={circle,fill,inner sep=0.9pt}, opt/.style={circle,draw,fill=white,inner sep=1.1pt}}
\begin{document}
\begin{tikzpicture}[x=1mm,y=1mm,line cap=round,line join=round,font=\scriptsize]
\node[inner sep=1pt] at (60,-13) {\tiny $N\cap N'=\mathfrak{z}=\pi_1$};
\node[inner sep=1pt] at (33,-27) {\scriptsize $N$};
\node[inner sep=1pt] at (72,-28) {\scriptsize $N'$};
\node[inner sep=1pt] at (52,-45) {\tiny $NN'=H$};
\node[inner sep=1pt] at (51,-70) {\scriptsize $G$};
\node[inner sep=1pt] at (17,-46) {\scriptsize $M'$};
\node[inner sep=1pt] at (83,-51) {\scriptsize $M$};
\draw[very thick] (34,-20.5) -- (42,-14);
\draw (64,-20) -- (67,-25);
\draw (36,-36) -- (42,-42);
\draw (64,-36) -- (57,-42);
\draw (51,-53) -- (51,-59);
\draw plot[smooth, tension=0.55] coordinates {(28,-34) (27,-48) (32,-60) (38.5,-67)};
\draw plot[smooth, tension=0.55] coordinates {(78.5,-39) (78,-51) (70,-62) (59.5,-70)};
\end{tikzpicture}
\end{document}
